\documentclass[11pt]{article}

\usepackage[margin=1in]{geometry}
\usepackage{amsmath, amssymb}
\usepackage{algorithm}
\usepackage{algpseudocode}
\usepackage{booktabs}
\usepackage{hyperref}

\title{Probabilistic graphical models}
\author{}
\date{}

\begin{document}
\maketitle

\section{Markov Networks and Conditional Random Fields (Lecture 6)}

This lecture covers the parameterization of Markov networks (MNs), their
log-linear form, conditional random fields (CRFs), and how MNs and Bayesian
networks (BNs) relate to and convert into one another.

\subsection{Independencies of a Markov network (recap)}

A trail $X_1 - \cdots - X_n$ is \emph{active} given $Z$ if none of the
$X_i$ (other than the endpoints) is in $Z$; unlike in a BN, activity of a
trail in an MN depends only on the graph structure, not on the direction of
edges or on colliders. This gives three families of independence statements
associated with an MN $\mathcal{H}$:

\begin{itemize}
  \item \textbf{Global Markov}: $\mathcal{I}(\mathcal{H}) =
    \{(\mathbf{X} \perp \mathbf{Y} \mid \mathbf{Z}) :
    \operatorname{sep}_{\mathcal{H}}(\mathbf{X}; \mathbf{Y} \mid \mathbf{Z})\}$,
    i.e. $\mathbf{X}$ and $\mathbf{Y}$ are independent given $\mathbf{Z}$
    whenever every trail between them is blocked by $\mathbf{Z}$.
  \item \textbf{Pairwise Markov}: $\mathcal{I}_p(\mathcal{H}) =
    \{(X \perp Y \mid \mathcal{X} - \{X,Y\}) : X\text{---}Y \notin \mathcal{H}\}$,
    any two non-adjacent nodes are independent given everything else.
  \item \textbf{Local Markov}: $\mathcal{I}_\ell(\mathcal{H}) =
    \{(X \perp \mathcal{X} - \{X\} - \mathrm{MB}_{\mathcal{H}}(X) \mid
    \mathrm{MB}_{\mathcal{H}}(X)) : X \in \mathcal{X}\}$, a node is
    independent of the rest of the graph given its Markov blanket (here,
    just its neighbors).
\end{itemize}

These three notions coincide for positive distributions; they are stated in
decreasing order of strength (global $\Rightarrow$ pairwise $\Rightarrow$
local in general, with equivalence under positivity).

\subsection{Gibbs distribution and alternative parameterizations}

The MN's graph being an I-map for $P$ is equivalent to $P$ factorizing over
the graph's cliques:

\[
P(\mathbf{X}) = \frac{1}{Z} \prod_{i=1}^{m} \phi_i(\mathbf{D}_i), \qquad
Z = \sum_{\mathbf{x} \in \mathrm{Val}(\mathbf{X})} \prod_{i=1}^m \phi_i(\mathbf{d}_i),
\]

where each $\mathbf{D}_i$ is (a subset of) a clique and $Z$ is the
\emph{partition function}, a normalizing constant. Writing $U(\mathbf{X}) =
\prod_i \phi_i(\mathbf{D}_i)$, $P(\mathbf{X}) = U(\mathbf{X})/Z$ is called a
\emph{Gibbs distribution}.

\subsubsection{Logarithmic / energy representation}

Since factors are non-negative, each can be written as
$\phi_i(\mathbf{D}_i) = e^{-\psi_i(\mathbf{D}_i)}$ for some
$\psi_i : \mathrm{Val}(\mathbf{D}_i) \to \mathbb{R}$. Substituting,

\[
P(\mathbf{X}) = \frac{1}{Z} \exp\Big(-\sum_i \psi_i(\mathbf{D}_i)\Big), \qquad
\ln P(\mathbf{X}) = -\ln Z - \sum_i \psi_i(\mathbf{D}_i).
\]

The quantity $\sum_i \psi_i(\mathbf{D}_i)$ is called the \emph{energy}:
lower energy corresponds to higher probability. A \emph{feature} is a
function $f : \mathrm{Val}(\mathbf{D}_i) \to \mathbb{R}$ that need not be
non-negative (a factor without the positivity requirement); a canonical
example is an \emph{indicator feature} that is $1$ for a particular
assignment and $0$ otherwise (e.g. ``these two superpixels have the same
color'').

\textbf{Worked example.} For the 4-cycle $A - B - C - D - A$ with pairwise
factors $\phi_1(A,B)$, $\phi_2(B,C)$, $\phi_3(C,D)$, $\phi_4(D,A)$, set
$\psi_i = -\ln \phi_i$ (so $\phi_1(a^0,b^0)=30 \Rightarrow \psi_1(a^0,b^0) =
-\ln 30 = -3.4$, etc.). Then
\[
P(A,B,C,D) \propto \exp\big[-\psi_1(A,B) - \psi_2(B,C) - \psi_3(C,D) - \psi_4(D,A)\big].
\]

\subsection{Log-linear Markov networks}

A \emph{log-linear} MN is specified by a set of features
$\mathcal{F} = \{f_1(\mathbf{D}_1), \dots, f_k(\mathbf{D}_k)\}$ (each with
scope $\mathbf{D}_i$, a complete subgraph) and weights $w_1, \dots, w_k$:

\[
P(X_1, \dots, X_n) = \frac{1}{Z} \exp\Big[-\sum_{i=1}^k w_i f_i(\mathbf{D}_i)\Big],
\qquad
\ln P(X_1,\dots,X_n) \propto \sum_{i=1}^k w_i f_i(\mathbf{D}_i).
\]

Different features may share a scope, and features/weights can be reused
across different factors (typically: features hand-designed, weights
learned from data). Log-linear form does not change the independence
assertions of the underlying model; it only changes how factors are
parameterized. Log-probability is linear in the weights, which is what
makes maximum-likelihood weight learning a (concave) convex optimization
problem.

\textbf{Example.} With $f_1(\text{Smoking}, \text{Cancer}) = \mathbb{1}[\lnot
\text{Smoking} \lor \text{Cancer}]$ and $w_1 = 1.5$, the feature/weight pair
encodes a soft implication ``smoking tends to cause cancer'' as one term in
$\sum_i w_i f_i$.

\subsubsection{Example: Ising model}

Binary spins $X_i \in \{+1,-1\}$ on a grid, with pairwise interaction
potentials $\phi_{i,j}(x_i,x_j) = w_{i,j} x_i x_j$ between neighbors and
single-node bias potentials $\phi_i(x_i) = u_i x_i$:

\[
P(\xi) = \frac{1}{Z}\exp\Big(-\sum_{i<j} w_{i,j} x_i x_j - \sum_i u_i x_i\Big).
\]

\subsubsection{Example: Boltzmann machine}

A (general) \emph{Boltzmann machine} is a fully-connected, pairwise MN
(Ising-like) over binary units $\mathbf{s} \in \{0,1\}^n$, split into
\emph{visible} units $\mathbf{v}$ (which carry the data) and \emph{hidden}
units $\mathbf{h}$ (which carry latent structure). Every pair of units --
visible-visible, hidden-hidden, and visible-hidden -- may be connected. Its
energy is a general quadratic form
\[
E(\mathbf{v},\mathbf{h}) = -\tfrac12\mathbf{v}^\top L \mathbf{v}
- \tfrac12\mathbf{h}^\top J \mathbf{h} - \mathbf{v}^\top W \mathbf{h}
- \mathbf{c}^\top \mathbf{v} - \mathbf{b}^\top \mathbf{h},
\]
with $P(\mathbf{v},\mathbf{h}) = \frac{1}{Z}\exp(-E(\mathbf{v},\mathbf{h}))$,
$Z = \sum_{\mathbf{v},\mathbf{h}} \exp(-E(\mathbf{v},\mathbf{h}))$. This is
just a log-linear MN whose features are all pairwise products $s_i s_j$ and
all singletons $s_i$, i.e.\ a fully general pairwise binary MRF. Because the
graph is a complete graph (every unit is a neighbor of every other unit),
none of the units are conditionally independent of each other given the
rest, so exact inference and even computing $Z$ are intractable in general
(the sum has $2^{|\mathbf{v}|+|\mathbf{h}|}$ terms); this is the price paid
for the extra representational power of hidden units and dense
connectivity.

\paragraph{Restricted Boltzmann machine (RBM).} Restricting the graph to be
\emph{bipartite} -- dropping all visible-visible edges ($L=0$) and all
hidden-hidden edges ($J=0$), keeping only visible-hidden edges $W$ -- gives
the RBM:
\[
E(\mathbf{v},\mathbf{h}) = -\mathbf{v}^\top W \mathbf{h} - \mathbf{c}^\top
\mathbf{v} - \mathbf{b}^\top \mathbf{h}, \qquad
P(\mathbf{v},\mathbf{h}) = \frac{1}{Z}\exp(-E(\mathbf{v},\mathbf{h})).
\]
The point of removing intra-layer edges is exactly the pairwise-Markov
property from Section~1.1: with no visible-visible or hidden-hidden edges,
each visible unit is separated from the other visible units by the hidden
layer (and vice versa), so
\[
P(\mathbf{h}\mid\mathbf{v}) = \prod_j P(h_j \mid \mathbf{v}), \qquad
P(\mathbf{v}\mid\mathbf{h}) = \prod_i P(v_i \mid \mathbf{h}),
\]
i.e.\ all hidden units are conditionally independent given the visible
units, and all visible units are conditionally independent given the
hidden units. Expanding the energy for a single unit shows each conditional
is a logistic/sigmoid unit exactly as in Section~2.2 (CRFs generalize
logistic regression):
\[
P(h_j = 1 \mid \mathbf{v}) = \operatorname{sigmoid}\Big(b_j + \sum_i v_i W_{ij}\Big),
\qquad
P(v_i = 1 \mid \mathbf{h}) = \operatorname{sigmoid}\Big(c_i + \sum_j W_{ij} h_j\Big).
\]
This is what makes RBMs practical: although $Z$ (and hence
$P(\mathbf{v},\mathbf{h})$ itself) is still intractable to compute exactly,
one can alternate sampling $\mathbf{h}\mid\mathbf{v}$ and
$\mathbf{v}\mid\mathbf{h}$ in \emph{block Gibbs sampling}, since each block
factorizes into independent coordinates. This alternating structure
underlies contrastive-divergence training, and stacking RBMs (feeding one
layer's hidden samples as the next RBM's visible data) is the basis of deep
belief networks.

\subsubsection{Metric MRFs}

A special case of pairwise MRFs used for labeling problems (image
segmentation, denoising): all $x_i$ take values in a common label space
$V$, and neighbors are encouraged to take similar labels via a smoothness
energy
\[
E(x_1,\dots,x_n) = \sum_i \epsilon_i(x_i) + \sum_{(i,j)\in\mathcal{E}} \epsilon_{i,j}(x_i,x_j),
\qquad P \propto \exp(-E),
\]
where $\epsilon_{i,j}$ is built from a distance/metric function $\mu$ on
labels, e.g. $\epsilon(x_i,x_j) = \min(c\|x_i-x_j\|_p,\ \mathrm{dist}_{\max})$
(a robust, truncated distance). Inference reduces to minimizing the energy,
equivalently maximizing the probability.

\section{Conditional Random Fields (CRFs)}

\subsection{Motivation}

Often we only care about $P(\mathbf{Y}\mid\mathbf{X})$ for target variables
$\mathbf{Y}$ given observed variables $\mathbf{X}$, not the full joint
$P(\mathbf{X},\mathbf{Y})$. Modeling $P(\mathbf{X},\mathbf{Y})$ directly
would force us to also model the (possibly complicated, high-dimensional,
or continuous) distribution over $\mathbf{X}$ itself. A CRF instead defines
an MN over $\{\mathbf{X}\cup\mathbf{Y}\}$ with factors
$\phi_1(\mathbf{D}_1),\dots,\phi_m(\mathbf{D}_m)$, but normalizes only over
$\mathbf{Y}$, for each fixed value of $\mathbf{X}$:

\[
P(\mathbf{Y}\mid\mathbf{X}) = \frac{1}{Z(\mathbf{X})}\tilde P(\mathbf{Y},\mathbf{X}),
\qquad
\tilde P(\mathbf{Y},\mathbf{X}) = \prod_{i=1}^m \phi_i(\mathbf{D}_i),
\qquad
Z(\mathbf{X}) = \sum_{\mathbf{Y}} \tilde P(\mathbf{Y},\mathbf{X}).
\]

This is parameterized exactly like a Gibbs distribution, but normalized
differently (per-$\mathbf{X}$ rather than globally). Consequences:
\begin{itemize}
  \item We never need to model dependencies among the $X$'s, which may be
    complex, poorly understood, continuous, or simply not of interest.
  \item We are free to use highly expressive, even overlapping/correlated
    features of $\mathbf{X}$ (e.g. many overlapping windows of text or
    image patches) without worrying about violating an independence
    assumption on $\mathbf{X}$ -- unlike, say, Naive Bayes, which assumes
    features are conditionally independent given the label and breaks down
    when features (e.g. neighboring pixels or words) are highly
    correlated.
  \item CRFs are widely used in NLP (sequence/entity labeling) and image
    segmentation.
\end{itemize}

\subsection{CRFs generalize logistic regression}

Logistic regression is the special case of a conditional MN with a single
binary query variable $Y$ and binary features $X_1,\dots,X_k$. With
pairwise potentials $\phi_i(X_i,Y) = \exp\{w_i \mathbb{1}[X_i=1,Y=1]\}$ and a
node potential $\phi_0(Y) = \exp\{w_0\mathbb{1}[Y=1]\}$,

\[
\tilde P(Y{=}1\mid x_1,\dots,x_k) = \exp\Big\{w_0+\sum_i w_i x_i\Big\}, \qquad
\tilde P(Y{=}0\mid \cdot) = 1,
\]
so that, after normalizing,
\[
P(Y{=}1\mid x_1,\dots,x_k) = \operatorname{sigmoid}\Big(w_0+\sum_{i=1}^k w_i x_i\Big),
\]
the familiar logistic regression model. This shows logistic regression is a
CRF with trivial (star-shaped) structure, and motivates CRFs as its
generalization to structured, correlated $\mathbf{Y}$.

\subsection{Applications}

\begin{itemize}
  \item \textbf{Image segmentation}: one node per superpixel $X_i$, edges
    between adjacent regions, a class label $S_i$ per superpixel. Because
    everything is conditioned on the (given) image, correlation among
    pixels does not need to be modeled; node/edge factors can use
    arbitrary image features (color histograms, texture, CNN outputs).
    Modern architectures (CRF-as-RNN, dense CRFs stacked on deep nets)
    embed CRF inference as a layer trained end-to-end with a CNN, rather
    than as a disconnected post-processing step.
  \item \textbf{Text/sequence labeling} (e.g. named entity recognition):
    target labels $Y_t$ (e.g. B-PER/I-PER/B-LOC/I-LOC/OTH) chained together
    with pairwise factors $\Phi_1(Y_t,Y_{t+1})$ enforcing label
    consistency, and per-position factors $\Phi_2(Y_t, X_1,\dots,X_T)$
    tying each label to (features of) the whole observed sequence.
  \item \textbf{Graph convolutional networks}: a CRF layer can be inserted
    between convolutional layers to encourage connected/similar nodes to
    have similar hidden representations, i.e. the CRF acts as a
    regularizer on the network's hidden features rather than on raw
    observations.
\end{itemize}

\section{Bayesian Networks vs. Markov Networks}

\subsection{Relative expressiveness}

\begin{itemize}
  \item \textbf{Directed (BN)}: captures inter-causal/``explaining away''
    reasoning; parameters (CPDs) are individually interpretable and can be
    elicited by hand; easy to generate data by ancestral sampling; large
    body of work on latent-variable models.
  \item \textbf{Undirected (MN)}: captures symmetric ``affinity''
    relationships; naturally represents cyclic dependency structures;
    parameters are less interpretable and usually learned from data;
    trickier (but tractable) parameter estimation; easy to add overlapping
    factors/features.
\end{itemize}

Neither family strictly contains the other: some independence patterns
(e.g. induced by a single common cause vs.\ a 4-cycle of pairwise
dependencies) can be captured exactly by a BN but not any MN, and vice
versa.

\subsection{BN $\to$ MN: moralization}

Every factor of a BN (i.e. every CPD, a node together with its parents)
must correspond to a subset of some clique of the target MN. \textbf{Moralize}
the graph: drop edge directions and add an (undirected) edge between any two
nodes that share a child (``marrying the parents''), so that every
node-plus-parents set becomes a clique. The moralized graph is a
\emph{minimal I-map} for the BN; it is a perfect map only if the original
DAG had no immoralities (v-structures) to begin with. Moralization can lose
independence information that was present in the original DAG (extra edges
= fewer independencies).

\subsection{MN $\to$ BN: triangulation}

Given an MN $\mathcal{H}$, build a minimal I-map BN by fixing an
elimination/variable ordering (each $X_i$ before its descendants) and, for
$i=1,\dots,n$, setting $\mathrm{Parents}(X_i)$ to the minimal subset
$\mathbf{S}\subseteq\{X_1,\dots,X_{i-1}\}$ such that $X_i \perp
(\{X_1,\dots,X_{i-1}\}\setminus\mathbf{S}) \mid \mathbf{S}$ in
$\mathcal{H}$. Applied along an ordering that visits the MN's cycle
structure, this typically introduces extra directed edges (to avoid
immoralities) -- the resulting BN must be \textbf{chordal} (every cycle of
length $\ge 4$ has a chord), since a v-structure-free DAG cannot contain
chordless long cycles. Equivalently, converting MN $\to$ BN requires
\textbf{triangulating} the MN's graph (adding edges until it is chordal);
different elimination orderings can produce very different numbers of
added edges, and the conversion can lose independence information.

\subsection{Perfect maps and chordal graphs}

Moralize (BN $\to$ MN) and triangulate (MN $\to$ BN) are, in general, lossy
in opposite directions. The distributions representable \emph{perfectly}
(as a P-map) in both families are exactly those whose MN is an
\textbf{undirected chordal graph} (a graph whose minimal cycles have length
$\le 3$). \textbf{Theorem}: if $\mathcal{H}$ is non-chordal, there is no BN
$G$ with $I(G) = I(\mathcal{H})$ -- because a minimal I-map BN for
$\mathcal{H}$ must itself be chordal, so representing a non-chordal MN
exactly as a BN is impossible without adding edges that destroy some of its
independence assertions.

\section{Factor Graphs}

A \textbf{factor graph} is a third graphical representation (alongside BNs
and MNs) that makes the factorization of a distribution explicit, rather
than only its independence structure. It is bipartite, with two kinds of
node:

\begin{itemize}
  \item \emph{variable nodes} (drawn as circles), one per random variable;
  \item \emph{factor nodes} (drawn as squares), one per factor $\phi$,
    connected to exactly the variables in its scope.
\end{itemize}

Unlike an MN graph -- which only records \emph{that} some factor exists
over each clique, collapsing all factors sharing that clique into one edge
set -- a factor graph distinguishes individual factors, including multiple
factors over the same or overlapping scopes (e.g.\ separate pairwise
factors $\phi_1,\phi_2,\phi_3$ plus higher-order factors $\phi_4$ over
triangle $\{X,Y,Z\}$, or extra unary factors $\phi_5,\phi_6,\phi_7$ on
individual variables). This makes factor graphs the natural representation
for deriving and running message-passing (sum-product / belief
propagation) inference algorithms, since messages are exchanged along the
explicit variable--factor edges.

\end{document}
